\documentclass[11pt]{article}
\usepackage[margin=1in]{geometry}
\usepackage{amsmath,amssymb,amsthm,mathtools}
\usepackage{enumitem}
\usepackage{array,booktabs}
\usepackage{microtype}
\usepackage[hidelinks]{hyperref}

\newcommand{\ALCIRCC}{\mathcal{ALCI}_{\mathrm{RCC5}}}
\newcommand{\comp}{\mathbin{;}}          % composition
\newcommand{\Comp}{\mathrm{comp}}
\newcommand{\code}[1]{\texttt{#1}}
\newcommand{\DR}{\mathsf{DR}}
\newcommand{\PO}{\mathsf{PO}}
\newcommand{\PP}{\mathsf{PP}}
\newcommand{\PPI}{\mathsf{PPI}}
\newcommand{\EQ}{\mathsf{EQ}}

\theoremstyle{plain}
\newtheorem{theorem}{Theorem}
\newtheorem{proposition}[theorem]{Proposition}
\newtheorem{corollary}[theorem]{Corollary}
\newtheorem{lemma}[theorem]{Lemma}
\theoremstyle{definition}
\newtheorem{definition}[theorem]{Definition}
\newtheorem{observation}[theorem]{Observation}
\newtheorem{remark}[theorem]{Remark}

\title{\bfseries The Width Barrier\\[2pt]
\large A Conditional Decidability Result for $\ALCIRCC$,\\
and the Unity of Its Two Open Directions}

\author{Claude (Fable~5, Anthropic)\\
\small prepared for Michael Wessel}
\date{July 14, 2026}

\begin{document}
\maketitle

\begin{abstract}
\noindent
Concept satisfiability for $\ALCIRCC$ --- the description logic
$\mathcal{ALCI}$ over the five RCC5 base relations under abstract
composition-table semantics --- has been open since Wessel
(2002/2003). This report records the current, honestly-delimited
state of one line of attack. We establish a \emph{conditional}
decidability theorem: $\ALCIRCC$ satisfiability is decidable provided a
single property holds --- \emph{bounded live width} (the obligation we
call F6) --- and we note that the soundness half of the underlying
reduction is machine-checked in Lean~4 with no unproved steps. We then
analyse F6 itself. Two facts about the RCC5 composition table, both
exhaustively verified, drive the analysis: (i) the only compositions
that \emph{determine} a relation (are singletons) all involve a
vertical ($\PP/\PPI$) argument, and $\PO$ is never a determined value;
(ii) consequently the ``determined'' structure a presentation may carry
for free (\emph{shadows}) is confined to the vertical axis, while the
costly \emph{live} structure lives on the horizontal ($\DR/\PO$) axis.
Building on these we show that the width barrier \emph{reduces to a
definability question}: F6 fails if and only if some $\ALCIRCC$ concept
can \emph{force} unboundedly large rigid horizontal configurations. The
substrate for such configurations is free (every $\DR/\PO$ network is
realizable), and horizontal chains \emph{block} exactly as vertical
ones do; so the entire barrier concentrates into whether the logic can
pin rigid horizontal coordinates. Finally we observe that this is
precisely the obstruction on the \emph{undecidability} side as well:
the reason no grid-encoding transfers to the abstract semantics
(Wessel 2002/2003; Lutz--Wolter 2006) is the same looseness. The
decidability keystone and the undecidability obstruction are one fact
seen from two sides. We do not resolve F6; we localize it, certify
everything around it, and place it on the knife's edge between the two
outcomes.
\end{abstract}

\section{Introduction}

The Region Connection Calculus fragment RCC5 partitions the
relationships between two non-empty spatial regions into five
jointly-exhaustive, pairwise-disjoint base relations. The description
logic $\ALCIRCC$ takes these five as its roles (inverse roles are
absorbed, since each relation's converse is again a base relation) and
interprets them over \emph{abstract} models: structures in which every
pair of distinct elements is labelled by exactly one base relation,
subject only to the RCC5 composition table and converse laws. Whether
concept satisfiability in this logic is decidable has been open for
over two decades \cite{wessel2003,lutzwolter2006}. No undecidability
reduction is known to apply to the abstract semantics, and no decision
procedure has been certified.

This report is not a claimed resolution. It is a \emph{structural
status report}: an account of where the difficulty now sits, written to
be handed to someone new. Its contributions are three.

\begin{enumerate}[leftmargin=1.6em]
\item A \textbf{conditional decidability theorem}
(\S\ref{sec:conditional}): $\ALCIRCC$ satisfiability is decidable if the
\emph{bounded live width} property F6 holds. The soundness direction of
the reduction --- that a finite, well-formed, composition-closed,
faithfully-labelled certificate yields a genuine RCC5 model of the
input concept --- is fully machine-checked in Lean~4.

\item A \textbf{reduction of F6 to a definability question}
(\S\S\ref{sec:anatomy}--\ref{sec:reduction}): using two exhaustively
verified facts about the composition table, we show that F6 fails
exactly when some concept can \emph{force} unboundedly large rigid
horizontal configurations, and that the combinatorial substrate for
such configurations is otherwise free.

\item A \textbf{unification of the two open directions}
(\S\ref{sec:knife}): the forcing question that F6 reduces to is the
same obstruction that blocks every attempt to prove the problem
\emph{undecidable}. Bounding it one way yields decidability; exploiting
it the other way would yield undecidability; and the reason neither has
happened is a single fact about the composition table.
\end{enumerate}

Throughout, we distinguish clearly between what is machine-checked, what
is exhaustively verified by finite computation, and what is argued.

\section{Preliminaries}\label{sec:prelim}

\subsection{RCC5 and its composition table}

The five base relations are $\DR$ (discrete --- disjoint interiors),
$\PO$ (partial overlap --- share a part, neither contained in the
other), $\PP$ (proper part), $\PPI$ (inverse proper part), and $\EQ$
(equal). Converse is the involution
$\overline{\DR}=\DR$, $\overline{\PO}=\PO$, $\overline{\EQ}=\EQ$,
$\overline{\PP}=\PPI$, $\overline{\PPI}=\PP$. The composition
$\Comp(R,S)$ collects every base relation that can hold between $x$ and
$z$ given $x\mathbin{R}y$ and $y\mathbin{S}z$; it is derivable from the
subset semantics on non-empty regions and is displayed in
Table~\ref{tab:comp}.

\begin{table}[h]
\centering
\renewcommand{\arraystretch}{1.25}
\setlength{\tabcolsep}{4pt}
\begin{tabular}{c|ccccc}
$\Comp$ & $\EQ$ & $\DR$ & $\PO$ & $\PP$ & $\PPI$\\\hline
$\EQ$  & $\{\EQ\}$ & $\{\DR\}$ & $\{\PO\}$ & $\{\PP\}$ & $\{\PPI\}$\\
$\DR$  & $\{\DR\}$ & $\{\EQ,\DR,\PO,\PP,\PPI\}$ & $\{\DR,\PO,\PP\}$ & $\{\DR,\PO,\PP\}$ & $\{\DR\}$\\
$\PO$  & $\{\PO\}$ & $\{\DR,\PO,\PPI\}$ & $\{\EQ,\DR,\PO,\PP,\PPI\}$ & $\{\PO,\PP\}$ & $\{\DR,\PO,\PPI\}$\\
$\PP$  & $\{\PP\}$ & $\{\DR\}$ & $\{\DR,\PO,\PP\}$ & $\{\PP\}$ & $\{\EQ,\DR,\PO,\PP,\PPI\}$\\
$\PPI$ & $\{\PPI\}$ & $\{\DR,\PO,\PPI\}$ & $\{\PO,\PPI\}$ & $\{\EQ,\PO,\PP,\PPI\}$ & $\{\PPI\}$\\
\end{tabular}
\caption{The RCC5 composition table, $\Comp(\text{row},\text{col})$.
Rows are the first argument, columns the second. Computed by exhaustive
enumeration over finite region domains.}
\label{tab:comp}
\end{table}

An \emph{atomic RCC5 network} on a set $V$ assigns to each ordered pair
of distinct elements exactly one base relation $N(u,v)$, satisfying
converse coherence $N(v,u)=\overline{N(u,v)}$ (R2) and composition
closure $N(u,w)\in\Comp(N(u,v),N(v,w))$ for all distinct $u,v,w$ (R3);
under \emph{strong-EQ} semantics $\EQ$ occurs only on the diagonal
(R1). We use the standard external input of the whole subject:

\begin{proposition}[Patchwork property; Renz--Nebel]\label{prop:patchwork}
A path-consistent atomic RCC5 constraint network is globally
realizable: it has a model.
\end{proposition}

\subsection{$\ALCIRCC$ syntax and semantics}

Concepts in negation normal form are generated by
\[
C,D \;::=\; \top \mid \bot \mid A \mid \neg A \mid C\sqcap D \mid C\sqcup D
\mid \exists R.C \mid \forall R.C,
\]
where $A$ is a concept name and $R\in\{\DR,\PO,\PP,\PPI\}$ (the reflexive
$\EQ$ is not used as a role between distinct elements). An
interpretation $\mathcal I=(\Delta,\rho,\cdot^{\mathcal I})$ has a
non-empty domain $\Delta$, an atomic RCC5 network $\rho$ on $\Delta$,
and concept-name extensions $A^{\mathcal I}\subseteq\Delta$. The role
reading is
$(\exists R.C)^{\mathcal I}=\{x\mid \exists y\,(\rho(x,y)=R \wedge y\in
C^{\mathcal I})\}$ and dually for $\forall$; because $\rho$ is atomic,
the extension of role $R$ is exactly the set of pairs labelled $R$.
$C_0$ is \emph{satisfiable} if $C_0^{\mathcal I}\neq\emptyset$ for some
such $\mathcal I$. This is the abstract composition-table semantics of
\cite{wessel2003}; it is what makes the problem different from the
topological logics of \cite{lutzwolter2006}, whose undecidability does
not transfer.

\section{The split-forest presentation and the two axes}\label{sec:axes}

The line of attack organizes any model along two axes. The
\emph{vertical} axis is nesting: $\PP$ and $\PPI$ form a forest, since
$\PP$ is a strict partial order ($\Comp(\PP,\PP)=\{\PP\}$, so $\PP$ is
transitive, and it is irreflexive under strong-EQ). The
\emph{horizontal} axis is $\DR$ and $\PO$: the relations between
elements that are not nested in one another. A model thus presents as a
forest of nested regions with horizontal cross-links between
incomparable ones; this is the \emph{split-forest presentation}. The
entire difficulty of the problem is captured by one asymmetry, which we
now make precise.

\section{Shadows and live width}\label{sec:width}

The composition table forces some relations and leaves others free. The
crucial instance is the pair
\begin{equation}\label{eq:downup}
\Comp(\PP,\DR)=\{\DR\}, \qquad
\Comp(\PPI,\DR)=\{\DR,\PO,\PPI\}.
\end{equation}
The first is a singleton: if $u\mathbin{\PP}w$ and $w\mathbin{\DR}v$
then $u\mathbin{\DR}v$ is \emph{forced} --- discreteness propagates
\emph{into parts}. The second is not: discreteness does not propagate
into wholes. This one-directional determinism is the source of both the
method's power and its difficulty.

\begin{definition}[Shadow / determined edge]\label{def:shadow}
In an atomic RCC5 network $N$ on $V$, an edge $\{u,v\}$ is
\emph{determined} (a \emph{shadow}) if its label is fixed by
path-consistency from $N$ restricted to $V\setminus\{u,v\}$ together
with the two elements' other edges. A sufficient local condition is the
existence of $w\in V$ with $|\Comp(N(u,w),N(w,v))|=1$. An edge that is
not determined is \emph{live}.
\end{definition}

A shadow costs nothing in a finite presentation: it is recoverable and
need not be independently stored or checked. Hence the size that
matters for a blueprint is not a node's degree but its number of live
neighbours.

\begin{definition}[Live width; the property F6]\label{def:width}
The \emph{live width} of a split-forest presentation is the supremum,
over its interface clusters, of the number of live occurrences that
must coexist in one cluster. Property \textbf{F6 (Bounded Live Width)}:
there is a computable $K$ such that every satisfiable concept $C_0$ has
a model whose split-forest presentation has live width at most
$K(|C_0|)$.
\end{definition}

\section{The conditional decidability theorem}\label{sec:conditional}

The reduction built over rounds of this project, and its soundness core
now certified in Lean~4, yields the following. We state it at the level
of certificates; the formal development is
\code{formal/Round19Transport.lean} (roughly 4{,}500 lines, zero
\code{sorry}, no added axioms).

\begin{theorem}[Conditional decidability]\label{thm:conditional}
If F6 holds, then $\ALCIRCC$ concept satisfiability is decidable.
\end{theorem}

\noindent\emph{Sketch, with the certified content marked.}
A \emph{certificate} is a finite \emph{catalogue} (core network,
attachment templates with a gluing normal form, and a template-indexed
steering function) together with a plan; unfolding a plan produces a
split-forest presentation. The development establishes, \emph{in the
Lean kernel}:
\begin{enumerate}[leftmargin=1.6em,itemsep=1pt]
\item \emph{Generation by construction.} For a structurally sound
catalogue and a valid plan, the unfolded certificate is well-formed and
\emph{faithful} (its pattern values agree with the recorded interface
rows) --- \code{build\_wellformed}, \code{build\_faithful}.
\item \emph{Catalogue soundness.} A finite composition check on the
catalogue alone (the \emph{S-condition}, independent of the plan length)
plus faithfulness implies the full per-unfolding closure condition ---
\code{scat\_scond} --- hence the unfolded frame is a proper atomic RCC5
network (composition-closed, converse-coherent, EQ-free):
\code{frame\_closed}, \code{catalogue\_soundness}.
\item \emph{Truth.} A Hintikka labelling of the frame with the target
concept at a root yields a genuine RCC5 model of that concept ---
\code{truth\_lemma}, \code{sat\_from\_hintikka}; and the induced frame
is verified to satisfy R1--R3 (\code{certInterp\_rcc5}).
\item \emph{Faithful abstraction.} Satisfiability and
Hintikka-realizability coincide (\code{satisfiable\_iff\_hintikkaP}): a
model always induces a Hintikka labelling, so the certificate loses
nothing.
\end{enumerate}
The decision procedure enumerates finite catalogues, checks the finite
S-condition and a Hintikka labelling, and accepts iff a valid one
carries $C_0$ at a root. Soundness (accept $\Rightarrow$ satisfiable) is
(1)--(3), certified. Completeness (satisfiable $\Rightarrow$ accept)
needs every satisfiable $C_0$ to admit a \emph{finite} catalogue; this
is the obligation \code{CompletenessObligation} in the development, and
it follows from F6 (which bounds the catalogue, giving a finite search
space and termination) together with a minor uniformization property
W2$'$ affecting only extraction faithfulness --- a failure of which can
only cause over-rejection, never unsoundness. \qed

\subsection{Soundness and completeness: what is certified}\label{sec:certstatus}

It is worth separating the two directions of the reduction and stating
exactly what the Lean development establishes for each, since the
completeness side is more subtle than a single ``machine-checked''
would suggest.

\emph{Soundness} --- a valid finite certificate yields a genuine RCC5
model --- is machine-checked in full (items (1)--(3) above:
\code{catalogue\_soundness}, \code{frame\_closed},
\code{sat\_from\_hintikka}).

\emph{Completeness} splits into two layers, and they have different
status.
\begin{itemize}[leftmargin=1.5em]
\item \textbf{Abstraction completeness (infinitary) --- certified.}
Allowing an arbitrary, possibly infinite Hintikka labelling,
completeness \emph{is} machine-checked: \code{satisfiable\_iff\_hintikkaP}
proves that $C_0$ is satisfiable \emph{if and only if} some RCC5 frame
carries a Hintikka labelling with $C_0$ at a point --- both directions,
no \code{sorry}. The non-trivial half, \code{model\_hintikkaP} (every
model induces a Hintikka labelling), is immediate once one observes that
the Hintikka clauses are exactly the recursion clauses of satisfaction
(\S\ref{sec:conditional}, item (4)). Thus the type-system abstraction is
\emph{provably faithful}: it captures satisfiability exactly and loses
nothing relative to models.
\item \textbf{Finitary completeness --- open, and equal to
$\mathrm{F6}\wedge\mathrm{W2}'$.} Requiring a \emph{finite} certificate,
as decidability demands, completeness is precisely the obligation
\code{CompletenessObligation}: every satisfiable $C_0$ admits a finite
catalogue, plan, and Hintikka labelling with $C_0$ at a root. This is
\emph{stated} in the development but deliberately \emph{not proved} and
\emph{not axiomatized}; it is the conjunction of F6 (bounded width, which
makes the catalogue finite and the search terminating) with W2$'$
(uniformization, which makes finite extraction faithful, and whose
failure can only over-reject).
\end{itemize}

\begin{center}
\renewcommand{\arraystretch}{1.2}
\begin{tabular}{@{}lll@{}}
\toprule
Direction & Lean witness & Status\\
\midrule
Soundness: certificate $\Rightarrow$ model
  & \code{sat\_from\_hintikka} & certified\\
Completeness (abstraction): model $\Leftrightarrow$ Hintikka
  & \code{satisfiable\_iff\_hintikkaP} & certified\\
Completeness (finitary): sat $\Rightarrow$ finite certificate
  & \code{CompletenessObligation} & open $=\mathrm{F6}\wedge\mathrm{W2}'$\\
Decision-grade reduction (shape only; premise oracle-inhabitable)
  & \code{BoundedDecider.decidable} & certified*\\
\bottomrule
\end{tabular}
\end{center}

\noindent In one line: the reduction is \emph{sound} and \emph{complete
modulo $\mathrm{F6}\wedge\mathrm{W2}'$} --- with soundness and the
infinitary abstraction-completeness both certified, and the sole
unformalized remainder being exactly $\mathrm{F6}\wedge\mathrm{W2}'$,
which the rest of this report analyses.

\begin{remark}
Theorem~\ref{thm:conditional} is a genuine reduction, not a
restatement: F6 is a single, precisely-posed property, and everything
around it is machine-checked. What is \emph{not} established is F6
itself (with the minor W2$'$). The rest of this report analyses it.
\end{remark}

\begin{remark}[Fourteenth review, GPT-5.5 cold]
The first cold review of the Lean development (July 2026) confirmed the
\emph{adequacy of the soundness pipeline} --- correct NNF syntax,
correct \code{sat} polarity and role reading, the correct strong-EQ
frame \code{RCC5Interp}, and correct composition/converse tables, with
no soundness defect and no vacuity. Its findings are formalization-level
and lie on the completeness interface: the certificate structures carry
higher-order \emph{function} fields rather than finite syntax;
\code{Satisfiable} fixes the carrier to the occurrence type rather than
arbitrary domains; and \code{CompletenessObligation} is an unbounded
existence rather than a bounded, decidable search statement. None is a
second mathematical conjecture beyond $\mathrm{F6}\wedge\mathrm{W2}'$;
each is the \emph{under-formalization} of the finitary step, and each is
repairable (finite codes with total evaluators; a countable-model
bridge; a decision-grade bounded obligation with an executable checker).
The review also confirmed one concrete over-restriction (the catalogue
check inspects unused diagonal template values; machine-verified,
completeness-side only). Thus the ``machine-checked surround'' of
Theorem~\ref{thm:conditional} is \emph{exact} for soundness and the
infinitary abstraction-completeness, while the finitary reduction still
requires the finite-code restatement above --- work orthogonal to, and
downstream of, F6.
\end{remark}

\begin{remark}[Decision-grade reduction, now formalized]
Responding to that review's central point, the decision-grade reduction
is now kernel-checked in the artifact. A \code{BoundedDecider} bundles a
computable bound, a decidable Boolean checker, and the two correctness
laws --- a passing code witnesses satisfiability, and every satisfiable
concept has a passing code below the bound --- and
\code{BoundedDecider.decidable} \emph{derives}
\code{Decidable (Satisfiable}~$C_0$\code{)} by a finite search over the
range below the bound. The conditional theorem is thereby genuinely
decision-grade: \code{Nonempty BoundedDecider} $\Rightarrow$
\code{Decidable}. What remains is to \emph{inhabit}
\code{BoundedDecider} --- precisely F6 (a computable width bound making
the search finite), W2$'$ (faithful finite extraction), and the
finite-coding of certificates --- but the passage from a bounded
decidable certificate to actual decidability is no longer a gap. (This is
the reduction the fourteenth reviewer's own sketch gestured at but could
not kernel-check for want of a Lean installation.)

Independently, the review's carrier finding (F2) is now closed:
\code{Interp} and \code{Satisfiable} are carrier-polymorphic, so
\code{Satisfiable} is \emph{arbitrary-domain} (reference) satisfiability
--- $\exists\,\alpha\;\exists\,I:\mathtt{Interp}\,\alpha\,\dots$ ---
with the occurrence-based pipeline as the $\alpha=\mathtt{Occ}$ instance.
The decision-grade theorem therefore concerns reference satisfiability,
and the carrier bridge (reducing an arbitrary-domain model to a
certificate) joins F6, W2$'$ and finite-coding on the single open
completeness side.
\end{remark}

\begin{remark}[Finite-coding, now formalized (F1)]
The review's remaining structural finding (F1) --- that \code{Cert},
\code{Template} and \code{Catalog} carry higher-order \emph{function}
fields rather than finite syntax --- is now addressed in the artifact.
The general fact is proved that a function is \emph{finite data on a
finite domain} (\code{fn\_finitely\_coded}: every $f$ has a finite
table reproducing it exactly on any finite set), and the finite code
types (\code{FinTemplate}, \code{FinCatalog}, all first-order data)
with \emph{total} decoders are supplied, exactly as the review
recommended. Faithfulness is proved for the first-order fields ---
\code{template\_net\_coded} and \code{coreNet\_coded} reconstruct
\code{net} and \code{coreNet} on their inspected port/index domains
--- and the higher-order steering field is handled by
\code{f\_factors\_through\_rows}: under the round-20
\code{f\_reads\_rows} clause, \code{f} depends only on the
\emph{finite} row-restrictions, so it too is finitely codable. (It is a
pleasing closure that the clause the soundness proof was forced to add
in round 20 is exactly what makes the steering function finite data.)
What remains is a finite-code bridge into the pipeline (see the
fifteenth-review correction below).
\end{remark}

\begin{remark}[Fifteenth review, GPT-5.5 cold --- corrections]
A cold review of rounds 26--28 (`papers/gpt5.5\_round28\_review/`)
returned \emph{gap, repairable but not cosmetic}, and it is correct; two
overclaims above are hereby withdrawn. \textbf{F2 is genuinely closed}
(round 27). \textbf{F3 is only partially reduced}: the reduction
\code{BoundedDecider.decidable} is correct, but its premise
\code{Nonempty BoundedDecider} is inhabitable by a \emph{classical
oracle} (\code{check}\,$C$\,$\_ :=$ \code{if Satisfiable}\,$C$
\code{then true else false}, verified to compile), so it proves
decidability from a classically vacuous premise rather than from a
finite syntactic certificate scheme. \textbf{F1 is not closed}: round 28
supplies correct \emph{local} tabulation lemmas and total decoders, but
the normative completeness witness remains higher-order (including the
labelling $\tau : \mathtt{Occ} \to \mathtt{List\;Concept}$), the
finite codes are not wired into the checker, and there is no global
inspected-domain coverage. The one \emph{immediately} repairable part
has been repaired: the non-vacuity witness
(\code{certK\_finitely\_codable}), which the review correctly called
\emph{degenerate} (empty domain, zero templates), now proves
\emph{total} agreement --- the decoded code reproduces \code{certK}'s
core net on all indices and its entire two-template list --- and a
companion \code{fin\_coding\_nontrivial} exhibits the tabulation
reconstructing a genuinely non-constant net (both constructive, no
axioms). The substantive F1 gap (a finite-code bridge \emph{into the
checker}, and a finite-code $\tau$) is untouched by this and remains
open. Toward closing F3, the decidable heart of the checker is now in the
artifact and machine-checked: \code{hintikkaB}, a Boolean
Hintikka-checker over the finite frame, is proved \emph{sound}
(\code{hintikkaB\_sound}: a passing check yields a genuine
\code{Hintikka} system), and \code{sat\_from\_hintikkaB} strengthens the
capstone so the labelling coherence --- the liveness heart, the one
acceptance test that is a finite computation yet had \emph{no}
\code{Decidable} instance as stated (the structural \code{Hintikka}
quantifies over the infinite \code{Nat}/\code{Concept} types) --- is
discharged by that finite Boolean test rather than \emph{assumed} (all
axiom-clean: \code{propext}/\code{Quot.sound} only, no \code{sorry}, no
choice). The residual, and it is the bulk of the finite-code bridge, is
the finite-table decidability of \code{Wellformed} and \code{SCond} ---
equivalently \code{CatOk}'s \code{net\_conv} and the core-net converse
as finite-table checks, \code{AttachOk}/\code{PlanOk}/\code{SCat} as
Boolean checks (each quantifies over finite lists), and the
by-construction \code{F\_reads} of a decoded \code{FinCatalog}. With
those the fixed non-oracular checker becomes total and F3 closes; it is
bounded, deliberate work, and honestly still open. \textbf{F4} also
remains. Accordingly the
claim ``decidable modulo exactly F6 $\wedge$ W2$'$'' is \emph{not yet
honest}: closing the reduction additionally requires a non-oracular
finite checker, a finite-code bridge for the certificate and the
labelling, global domain coverage, and the F4 repair. What is genuinely
established is narrower: soundness and infinitary abstraction-completeness
(certified), the reference-semantics target (F2, certified), and a
correct finite-search reduction \emph{shape} awaiting an honest checker.
The open mathematics F6 and W2$'$ are unchanged and remain the deep
obstruction; the interface repairs above are finishable formalization,
not new mathematics.
\end{remark}

\begin{remark}[Round-29: the non-oracular finite checker --- F3 closed, F1 wired]
The finishable formalization above has now been finished, and
machine-checked (all \code{propext}/\code{Quot.sound} only; no
\code{sorry}, no choice; \code{native\_decide} only in the witness). The
artifact now carries a \emph{fixed}, non-oracular Boolean checker
\code{finAcceptB} of a \emph{first-order finite} code --- a coded
catalogue (\code{FinCatalog}), a plan, a labelling \emph{table}, and a
root --- that mentions neither \code{Satisfiable} nor any oracle, only
finite table/decision checks and the Boolean Hintikka checker
\code{hintikkaB}. Its soundness is a theorem (\code{finAccept\_sound}: a
passing check yields an RCC5 model, via the whole pipeline
--- converse tables give net-coherence, the bridge decoder gives
\code{F\_reads} by construction, the bounded fields are decided,
\code{catalogue\_soundness} and \code{build\_wellformed} give
\code{SCond}/\code{Wellformed}, and \code{sat\_from\_hintikkaB} gives the
model). The decision-grade reduction \code{decidableSat\_of\_finScheme}
then derives \code{Decidable (Satisfiable}~$C_0$\code{)} from a fixed
enumeration and a completeness premise \emph{about the fixed checker}.
Because the checker is fixed and forces exhibiting a real finite
certificate, that premise is \emph{not} oracle-inhabitable (a
\code{Satisfiable} proof yields a possibly-infinite model, not a finite
code) --- unlike round-26's free-field \code{BoundedDecider}, which it
supersedes. Non-vacuity is witnessed (\code{finAccept\_nonvacuous}:
\code{finAcceptB} returns \code{true} on \code{certK} coded as a
\code{FinCatalog}, with \code{atom 0} at the root). So the fifteenth
review's F3 (oracle-vacuity) is \emph{closed}, and F1's finite codes ---
catalogue \emph{and} labelling --- are now \emph{wired into} a total
checker. What remains is exactly the completeness premise itself: a
\emph{computable, complete} enumeration, which is precisely the open
mathematics F6 $\wedge$ W2$'$. This does \emph{not} prove decidability;
it makes ``decidable modulo a complete finite-certificate enumeration''
an honest, kernel-checked reduction, with the residual premise pinned to
F6 $\wedge$ W2$'$. Status unchanged: strongly supported, not certified.
\end{remark}

\begin{remark}[Round-30: F4 fixed --- the last interface finding closed]
The one interface finding still outstanding after round-29 --- F4, the
\code{SCat.net\_r3} over-restriction --- is now repaired and
machine-checked. The catalogue check enforced its triangle law for
\emph{all} member-port pairs, including the degenerate case where a port
\emph{is} the fresh port $\mathtt{inr(inr\ }j)$, where it read the
template's junk diagonal self-value and could spuriously reject a valid
certificate (\code{PP}$\notin\Comp(\code{PP},\code{conv DR})=\{\code{DR}\}$;
machine-verified). It now guards $p\neq\mathtt{inr(inr\ }j)$ and
$q\neq\mathtt{inr(inr\ }j)$, excluding exactly those junk triples. The
repair is \emph{sound} because the sole consumer of \code{net\_r3} ---
\code{fresh\_tri} --- never needs the degenerate case: its member ports
map to real occurrences $x,y\neq\mathtt{born}\ i\ j$, so it supplies the
two guards from its own hypotheses, and the soundness chain
(\code{fresh\_tri}$\to$\code{scat\_scond}$\to$\code{catalogue\_soundness})
is unchanged. So all four of the fifteenth review's interface findings
(F1--F4) are now addressed, and the only open items are the mathematics
F6 $\wedge$ W2$'$. Status unchanged: strongly supported, not certified.
\end{remark}

\section{The anatomy of forcing}\label{sec:anatomy}

We now isolate what the composition table can and cannot \emph{force}.
Call a composition $\Comp(R,S)$ \emph{determining} if it is a singleton.

\begin{proposition}[Determination is vertical]\label{prop:singletons}
Among the sixteen compositions of proper (non-$\EQ$) relations, exactly
four are determining:
\[
\Comp(\DR,\PPI)=\{\DR\},\quad
\Comp(\PP,\DR)=\{\DR\},\quad
\Comp(\PP,\PP)=\{\PP\},\quad
\Comp(\PPI,\PPI)=\{\PPI\}.
\]
Every one of them has a vertical ($\PP$ or $\PPI$) argument. No
composition with both arguments horizontal ($\in\{\DR,\PO\}$) is
determining, and $\PO$ is never the value of a determining composition.
\end{proposition}

\begin{proof}
Exhaustive inspection of Table~\ref{tab:comp}; verified by finite
computation. The horizontal block $\{\DR,\PO\}\times\{\DR,\PO\}$ gives
$\Comp(\DR,\DR)$ and $\Comp(\PO,\PO)$ (all five relations) and
$\Comp(\DR,\PO)=\Comp(\PO,\DR)$'s companions of size three; none is a
singleton. The four singletons are as listed, each with a $\PP/\PPI$
leg. Scanning the table, no cell equals $\{\PO\}$.
\end{proof}

Proposition~\ref{prop:singletons} explains a phenomenon observed
experimentally: attempts to force a region to acquire unboundedly many
neighbours succeed only in producing shadows. Any \emph{forced} edge
arises from a determining composition (or an intersection of
constraints that pins a value), and by
Proposition~\ref{prop:singletons} the only fully determining
compositions run through the vertical axis --- so forced structure is
shadow structure, confined to nesting. Concretely, a construction in
which a distinguished ``spy'' region is forced discrete from every
level of an ascending tower makes the spy's degree grow without bound,
yet each such edge is a shadow (forced by $\Comp(\PP,\DR)=\{\DR\}$
through the tower); the \emph{live} width stays bounded.

There is, however, a subtlety that keeps the barrier open. Determination
and \emph{distinctness} are different.

\begin{proposition}[Distinctness can be horizontal]\label{prop:distinct}
Two same-typed occurrences are mergeable (identifiable, set to $\EQ$)
unless some constraint excludes $\EQ$ between them. A shared horizontal
anchor excludes $\EQ$: if $u\mathbin{\DR}w$ and $w\mathbin{\PO}v$ then
$u,v$ satisfy $\Comp(\DR,\PO)=\{\DR,\PO,\PP\}$, which excludes $\EQ$
(forcing $u\neq v$) yet is not a singleton (the edge $\{u,v\}$ remains
live).
\end{proposition}

\begin{proof}
From Table~\ref{tab:comp}, $\Comp(\DR,\PO)=\{\DR,\PO,\PP\}\not\ni\EQ$
and $|\{\DR,\PO,\PP\}|=3$.
\end{proof}

Thus \emph{determination} is strictly vertical
(Prop.~\ref{prop:singletons}) but \emph{distinctness} extends to the
horizontal axis (Prop.~\ref{prop:distinct}). The two notions come
apart, and it is precisely in the gap between them that F6 lives: a
live horizontal crowd --- many occurrences, pairwise distinct by
horizontal anchors, with live (undetermined) mutual relations --- is
not ruled out by Proposition~\ref{prop:singletons}.

\section{The width barrier reduces to horizontal forcing}\label{sec:reduction}

We now follow the live-crowd possibility to its floor. Three facts,
each finite-checkable, combine into a reduction.

\begin{proposition}[The substrate is free]\label{prop:substrate}
Every $\{\DR,\PO\}$-labelled network is path-consistent, hence
realizable by Proposition~\ref{prop:patchwork}. Moreover, for every $n$
there is such a network on $n$ nodes in which no two nodes are
mergeable (all rows distinct): an irreducible live horizontal crowd of
size $n$.
\end{proposition}

\begin{proof}
For $R,S\in\{\DR,\PO\}$ one has $\{\DR,\PO\}\subseteq\Comp(R,S)$ (rows
$\DR,\PO$ of Table~\ref{tab:comp}), so any $\{\DR,\PO\}$-labelling meets
every composition constraint: path-consistency is immediate.
Realizability follows from Proposition~\ref{prop:patchwork}. For
irreducibility, a $\{\DR,\PO\}$ network is a graph (edge $=\PO$,
non-edge $=\DR$); a graph with pairwise-distinct neighbourhood rows has
no two mergeable nodes, and such graphs exist at every size.
\end{proof}

\begin{proposition}[Horizontal chains block]\label{prop:block}
The concept $\exists\PO.\top \sqcap \forall\PO.\exists\PO.\top$, which
demands an endless horizontal successor chain, is satisfiable by a
finite model. More generally, any forced $\PO$-chain admits a finite
looping model.
\end{proposition}

\begin{proof}
$\Comp(\PO,\PO)$ is the whole algebra
(Table~\ref{tab:comp}), and in particular contains $\EQ$; hence a
$\PO$-chain may close a loop ($w_i=w_j$ is composition-consistent). With
finitely many concept types, any chain repeats a type, and a repeated
type with mergeable surroundings closes the loop, exactly as in the
vertical blocking (\S\ref{sec:axes}). A two- or three-element cyclic
model satisfies the displayed concept.
\end{proof}

\begin{proposition}[No non-EQ horizontal recurrence --- path form]\label{prop:pathauto}
Form the composition-path automaton whose states are the reachable
relation sets $\mathrm{Rel}(w)$, with $\mathrm{Rel}(\varepsilon)=\{\EQ\}$
and $\mathrm{Rel}(w\cdot a)=\bigcup_{r\in\mathrm{Rel}(w)}\Comp(r,a)$,
restricted to states not containing $\EQ$. Every cycle in this
restricted automaton uses only vertical labels ($\PP$ or $\PPI$); no
cycle uses a horizontal label ($\DR$ or $\PO$).
\end{proposition}

\begin{proof}
Exhaustive finite computation: there are $11$ reachable relation sets,
$8$ without $\EQ$, and the six cyclic edges among them are the
self-loops $\{\PP\},\{\PO,\PP\},\{\DR,\PO,\PP\}$ under $\PP$ and
$\{\PPI\},\{\DR\},\{\DR,\PO,\PPI\}$ under $\PPI$; none is horizontal.
\end{proof}

Proposition~\ref{prop:pathauto} strengthens
Proposition~\ref{prop:singletons} from single compositions to whole
paths, and sharpens Proposition~\ref{prop:block}: \emph{any} support
recurrence that never admits an equality fold is eventually vertical,
hence absorbed by vertical cyclic unfolding, not by a horizontal crowd.
It removes a class of would-be counterexamples but does not close F6,
which is a statement about \emph{width} rather than about paths. (This
proposition was contributed by the fourteenth review and independently
verified.)

Propositions~\ref{prop:substrate} and~\ref{prop:block} isolate the
issue. A large live horizontal crowd is trivial to \emph{exhibit}
(Prop.~\ref{prop:substrate}) but cannot be produced merely by
\emph{demanding} more occurrences, because horizontal chains fold back
on themselves (Prop.~\ref{prop:block}). To prevent the fold one would
have to make the crowd members pairwise \emph{non-mergeable in every
model} --- to give each a permanent address, a rigid horizontal
coordinate. By Proposition~\ref{prop:singletons} the composition table
offers no lever to pin a horizontal coordinate: it cannot determine a
horizontal relation, and $\PO$ is never forced. We arrive at the central
structural statement.

\begin{observation}[The width barrier is a forcing question]\label{obs:reduction}
Bounded live width (F6) fails if and only if some $\ALCIRCC$ concept
\emph{forces} arbitrarily large \emph{rigid} horizontal configurations
--- unboundedly many occurrences that must coexist at one interface and
that no model may merge. Equivalently:
\begin{center}
\emph{F6 counterexample} $\iff$ \emph{a concept forcing an unbounded
rigid $\{\DR,\PO\}$ graph},\\
\emph{F6 proof} $\iff$ \emph{no concept forces such a graph.}
\end{center}
The combinatorial substrate is free (Prop.~\ref{prop:substrate}); the
mere generation of occurrences blocks (Prop.~\ref{prop:block}); so the
whole barrier concentrates into whether the logic can force rigid
horizontal coordination, a \emph{definability} question about the
expressive power of $\ALCIRCC$.
\end{observation}

Observation~\ref{obs:reduction} is a reduction, not a theorem in the
formal sense: the forward direction (unbounded width $\Rightarrow$
forced irreducible crowds) is the definition of live width unwound; the
content is that the converse obstruction is a pure forcing question,
because everything \emph{except} forcing --- realizability and
non-repetition --- has been discharged by
Propositions~\ref{prop:substrate} and~\ref{prop:block}. We state it as
an observation to mark that ``forces'' is used in its intuitive
model-theoretic sense; making it a formal equivalence is itself part of
the open problem.

\section{One fact, two faces}\label{sec:knife}

The forcing question of Observation~\ref{obs:reduction} is not new to
the subject. It is the exact obstruction on the \emph{undecidability}
side.

To prove a logic undecidable one typically forces it to encode an
unbounded, rigidly-coordinated grid on which a Turing computation (or a
tiling) is written. For RCC8 with \emph{topological} semantics this
succeeds \cite{lutzwolter2006}. For the \emph{abstract}
composition-table semantics of $\ALCIRCC$ (and $\mathcal{ALCI}_{\mathrm
{RCC8}}$) it does not, and the reason was identified by
\cite{wessel2003}: the composition table cannot force the
\emph{coincidence condition} --- that ``east-of-north'' and
``north-of-east'' name the same cell --- because it cannot pin the
relevant relations to single values. That is rigid horizontal
coordination, and its unavailability is exactly
Proposition~\ref{prop:singletons}: no horizontal determination, $\PO$
never forced.

\begin{observation}[The knife's edge]\label{obs:knife}
The bounded-width property F6 (decidability side) and the
non-forcibility of a coordinate grid (undecidability side) are the same
fact about $\ALCIRCC$: \emph{the abstract composition table is too
loose to pin rigid horizontal structure}. If this looseness can be
turned into a bound, one obtains decidability
(Theorem~\ref{thm:conditional}); if it could be circumvented and a grid
forced, one would obtain undecidability. The two open directions are one
question seen from two sides, and the reason neither has been settled in
two decades is that the same fact must be turned one way or the other.
\end{observation}

\section{The cold-attack sharpening: identity-selector minors}\label{sec:sharpen}

A cold research attack on this report (GPT-5.5, on the F6/W2$'$ packet
\code{papers/cold\_review\_f6\_w2prime}) did not settle F6 --- as
expected --- but it \emph{sharpened} the target and eliminated several
would-be shortcuts. The machine-checkable witnesses below were
reproduced locally (\code{wp39}--\code{wp44}); the accompanying
theorem-level claims are the attacker's and are \emph{not} yet
independently verified.

\paragraph{``Rigid horizontal crowd'' is too broad.}
Observation~\ref{obs:reduction} asked whether a concept can force
unboundedly many pairwise non-mergeable horizontal occurrences. That
condition is satisfied by constructions that are nonetheless harmless:
a shell half-graph ($y_i\mathbin{\DR}p_j$ for $j\le i$,
$y_i\mathbin{\PO}p_j$ for $j>i$) forces an unbounded raw horizontal
crowd, yet all shells share the same relevant row at the recursive
interface, so a finite row-keyed catalogue represents the whole family
(\code{wp41}/\code{wp42}); a prefix-monotone ``pumped'' staircase is
finite-state with one moving inherited slot (\code{wp43}); and a
single-chain PO-gap ($\exists\PO.A$ in one phase, $\forall\PO.\neg A$ in
another) has a \emph{bounded live window} --- its PO-events are singleton
shadows except three local edges, so at most two live edges meet each
chain point (\code{wp40}). Raw non-mergeability is thus the wrong
invariant; the certificate-relevant one is unbounded row/profile
inequivalence that survives bounded refinement.

\paragraph{The sharp target.}
The precise obstruction is an \emph{$M_n$ identity-selector minor}: pairs
$(L_i,U_i)_{i<n}$ with selectors $(A_i)_{i<n}$ and a fresh splitter $b$
such that $A_i$ separates $L_i,U_i$ (i.e.\ $\rho(L_i,A_i)\neq
\rho(U_i,A_i)$) while $A_j$ does \emph{not} for $j\neq i$, and the
splitter forces $L_i$ and $U_i$ apart ($L_i\mathbin{\DR}b$,
$U_i\mathbin{\PPI}b$). The selector incidence matrix is the $n\times n$
identity; the minimum \emph{separator-cover} is $n$. Half-graphs and
prefix pumps have separator-cover $O(1)$ (one carried memory slot);
identity minors have cover $n$ by design. The reduction sharpens to:
\begin{center}
\emph{F6 holds} $\iff$ \emph{no finite $\ALCIRCC$ concept forces
unbounded $M_n$ identity-selector minors};\\
\emph{F6 fails} $\iff$ \emph{some concept forces them for all $n$.}
\end{center}

\paragraph{An algebra-only proof of F6 is refuted.}
One tempting route --- bound the live width by compressing
path-completion support trees --- is blocked. For every $n$ there is a
finite, composition-closed RCC5 \emph{set} network realizing an $M_n$
minor with separator-cover exactly $n$, while each individual support
tree has height one (\code{wp44}, verified through $n=8$). So bounded
support depth does \emph{not} imply bounded separator cover. Crucially
this family is a bare RCC5 network, \emph{not} forced by any concept, so
it refutes a proof \emph{strategy}, not F6. Any F6 proof must use that
the bad glue steps are $\ALCIRCC$-\emph{unforceable}, not merely that
they are RCC5-consistent. The recommended positive route is a
\emph{selector-sharing lemma} over bounded \emph{tuple} profiles (not
unary types) --- if many same-profile selector gadgets must share or
collapse, unbounded identity minors cannot be forced --- which is a
back-and-forth/bisimulation argument, not a locality one (horizontal
degree is unbounded, so Gaifman locality alone does not apply).

\paragraph{W2$'$ folds into F6.}
The uniformization property W2$'$ is not an independent obstacle. Its
coarse form is \emph{false}: a five-point witness (\code{wp39}) has two
old occurrences $y,z$ of identical coarse state (type, separator row,
safe domain to the fresh $b$) whose every joint completion forces
$y\mathbin{\DR}b$ but $z\mathbin{\PPI}b$ --- separated only by their
relation to a third occurrence $a$ \emph{outside} the separator. The
failure is completeness-side (over-rejection, never unsoundness), and it
is repaired by refining the interface state with the row to the bounded
set of \emph{live} old anchors; since F6 bounds that set, the refinement
is bounded \emph{by F6}. Thus the honest conditional is F6 plus a
\emph{live-neighbourhood uniformization lemma} whose size is controlled
by F6 itself --- one open problem, not two.

\paragraph{A decidable fragment.}
Finally, one unconditional fragment falls out. In the two-tier quotient
analysis $\PO$ is the only relation whose recurrence needs an all-phase
safety condition; a concept with \emph{no $\forall\PO$ subformula} makes
that condition vacuous, so it admits a finite two-tier quotient and a
computable live-width bound $K(C)=(1+m)^2\,2^{2^m}$ with
$m=|\mathrm{cl}(C)|$.
Hence the \emph{$\forall\PO$-free fragment} of $\ALCIRCC$ (retaining
$\exists\PO$, $\forall\DR$, and all vertical modalities) is decidable
--- a Target-B win independent of full F6, resting on the two-tier
manuscript's stabilization and forcing/absorption lemmas
(theorem-level, not yet certified).

\section{Discussion}\label{sec:discussion}

\paragraph{Why the combinatorics did not fit.}
Attempts to attack the width bound with Ramsey-type combinatorics were
unsuccessful, and Observation~\ref{obs:reduction} explains why by type
mismatch. Ramsey theory concerns structure that is \emph{unavoidable}
in a large fixed object; F6 concerns structure that a formula can
\emph{force}. These differ in quantifier order: ``every large graph
contains $X$'' versus ``some formula makes every model contain $X$.''
The tools that fit the latter are model-theoretic: locality
(Gaifman/Hanf), Ehrenfeucht--Fra\"iss\'e games, the composition method,
and the finite/bounded-model-property machinery --- together with the
grid-encodability analysis in which
\cite{wessel2003,lutzwolter2006} already sit. This is a different
toolbox from the model-\emph{construction} techniques used to build the
certificate layer, and it is the one Observation~\ref{obs:knife}
indicates. We do not claim it succeeds; we claim it is the correct room
to enter.

\paragraph{What is and is not established.}
Certified (Lean, no gaps): the soundness of the certificate-to-model
pipeline and the faithfulness of the Hintikka abstraction
(\S\ref{sec:conditional}). Exhaustively verified (finite computation):
Table~\ref{tab:comp} and Propositions~\ref{prop:singletons},
\ref{prop:distinct}, \ref{prop:substrate}, \ref{prop:block}. Argued
(structural, intuitive ``forces''): Observations~\ref{obs:reduction}
and~\ref{obs:knife}. Open: F6 itself, hence decidability. The honest
one-line status is unchanged from the project's standing label:
\emph{strongly supported, not certified}.

\paragraph{Value independent of F6.}
Theorem~\ref{thm:conditional} is a conditional result that stands
regardless of F6's fate, with its soundness half machine-certified.
Observation~\ref{obs:knife} is, to our knowledge, a new unification: it
identifies the decidability keystone and the undecidability obstruction
as one fact. Together they turn a formless two-decade question into a
single, well-posed, two-sided lemma with a certified surrounding and a
named toolbox. If the problem is ever settled we expect the resolution
to begin from a map of this shape.

\section{Conclusion}

$\ALCIRCC$ satisfiability is decidable if the live width of its models
is bounded (Theorem~\ref{thm:conditional}), and the soundness of that
reduction is machine-checked. Whether the width is bounded reduces to
whether the logic can force rigid horizontal coordination
(Observation~\ref{obs:reduction}), which is exactly the obstruction
that blocks proving the problem undecidable
(Observation~\ref{obs:knife}). The composition table's inability to pin
rigid horizontal structure is the single fact underneath both open
directions. We have not turned that fact into a bound; we have shown
that doing so, in either direction, is the whole of what remains.

\begin{thebibliography}{9}
\bibitem{wessel2003} M. Wessel.
\emph{On spatial reasoning with description logics --- position paper}
and technical report (report7), 2002/2003. Introduces
$\mathcal{ALCI}_{\mathrm{RCC}}$ under abstract composition-table
semantics and identifies the coincidence-condition obstruction to grid
encoding.
\bibitem{lutzwolter2006} C. Lutz and F. Wolter.
\emph{Modal Logics of Topological Relations}. Logical Methods in
Computer Science, 2006. Undecidability for topological RCC logics; the
reductions do not transfer to the abstract semantics.
\bibitem{renznebel} J. Renz and B. Nebel.
\emph{On the complexity of qualitative spatial reasoning: a maximal
tractable fragment of the region connection calculus}. Artificial
Intelligence, 1999. Source of the patchwork property.
\end{thebibliography}

\vfill
\noindent\footnotesize
Companion artifacts: the machine-checked development
\code{formal/Round19Transport.lean}; the width-attack probes
\code{verification/python/wp35\_f6\_width\_attack.py},
\code{wp36\_determination\_vs\_distinctness.py},
\code{wp37\_f6\_reduces\_to\_forcing.py}; and the plain-language
companion \code{WHY\_ITS\_HARD.md}. This report is deliberately at the
boundary of formal and expository; where it and the machine-checked
artifacts disagree, the artifacts win.
\end{document}
